\documentclass[a4paper]{article}
\usepackage{amsfonts}
\usepackage{amsmath}
\usepackage{amssymb}
\usepackage{graphicx}     

\begin{document}
  
\noindent \large Auteur: Noran Ben Said \\
\noindent \large Studierichting: Informatica\\
\noindent \large Oefeningenreeks 5A nr. 7.4 \\
\medskip

\normalsize

$ \displaystyle \int_0^{\pi} \cos^2 x \sin^3 x \, dx$\\

$ = \displaystyle \int_0^{\pi} \cos^2 x \sin^2 x \sin x \, dx$\\

$ = \displaystyle \int_0^{\pi} \cos^2 x (1-\cos^2 x)\sin x \, dx$\\

Stel $ t = \cos x \Rightarrow dt = -\sin x \, dx$\\

Nieuwe grenzen:\\

Als $x=0$, dan $t=\cos 0=1$\\

Als $x=\pi$, dan $t=\cos \pi=-1$\\

$ = \displaystyle \int_1^{-1} t^2(1-t^2)(-dt)$\\

$ = \displaystyle \int_{-1}^{1} t^2(1-t^2)\,dt$\\

$ = \displaystyle \int_{-1}^{1} (t^2-t^4)\,dt$\\

$ = \displaystyle \int_{-1}^{1} t^2\,dt - \int_{-1}^{1} t^4\,dt$\\

$ = \left[\dfrac{t^3}{3}\right]_{-1}^{1} - \left[\dfrac{t^5}{5}\right]_{-1}^{1}$\\

$ = \left(\dfrac{1}{3}-\left(-\dfrac{1}{3}\right)\right) - \left(\dfrac{1}{5}-\left(-\dfrac{1}{5}\right)\right)$\\

$ = \dfrac{2}{3} - \dfrac{2}{5} = \dfrac{4}{15}$\\


\begin{thebibliography}{999}
%\bibitem{a} Referentie
\end{thebibliography}
\end{document} 
